\section{Algorithm Explanation}
\vspace{1cm}
Generative models offer a powerful approach for compressed sensing by leveraging the structure of the data distribution.
\\
Independently from the above mentioned models, a generative model is defined by a deterministic function \( G: \mathbb{R}^k \to \mathbb{R}^n \) and a distribution \( P_Z \) over \( z \in \mathbb{R}^k \). The model maps a low-dimensional latent vector \( z \) (decided by us as the latent dimension) to the high-dimensional sample space (the final generated image). The compressed sensing task involves finding a vector \( z \) in the latent space such that the corresponding generated sample \( G(z) \) matches the observed measurements \( y \).
\\
The measurement matrix \( A \) is drawn at random, with entries sampled independently from a Gaussian distribution \( N(0, 1/m) \). With this scaling the operator preserves norms in expectation, since \( \mathbb{E}\|Ax\|^2 = \|x\|^2 \) for every \( x \), and the magnitude of the measurement error remains comparable across different values of \( m \).
\\
The objective function to be minimized is:\\
\begin{equation}
\text{loss}(z) = \| AG(z) - y \|^2.
\end{equation}
The objective actually used in \cite{bora2017} adds a regularization term to this expression:
\begin{equation}
\text{loss}(z) = \| AG(z) - y \|^2 + \lambda \|z\|^2 .
\end{equation}
Since both the VAE and the DCGAN impose an isotropic Gaussian prior on \( z \), the quantity \( \|z\|^2 \) is proportional to the negative log-likelihood of the latent vector under that prior, so the term keeps the search inside the region the generator was trained to cover. The authors report \( \lambda = 0.1 \) as the best value on MNIST and we adopt it, applying it to every generative prior for consistency. Results obtained with \( \lambda = 0 \) are reported alongside, as the reference paper does.
\\
Optimization is performed using gradient descent on \( z \) with the Adam algorithm \cite{kingma2015adam} for optimizing descent. If \( G \) is differentiable, gradients of the loss with respect to \( z \) can be computed using back-propagation. The process involves iteratively updating \( z \) to minimize the measurement error. Upon convergence to \( \hat{z} \), the reconstruction of \( x^* \) is given by \( G(\hat{z}) \).
\\\\
The algorithm can be summarized as follows:
\begin{enumerate}
    \item Initialize \( z \) by sampling from the distribution \( P_Z \), which for both models is the standard normal \( N(0, I) \) used during training.
    \item Compute the loss: \( \text{loss}(z) = \| AG(z) - y \|^2 \).
    \item Use gradient descent to update \( z \) to minimize the loss.
    \item Repeat until convergence.
    \item The reconstructed signal is \( \hat{x} = G(\hat{z}) \).
\end{enumerate}
This method leverages the data prior encoded in the generative model to achieve high-quality reconstructions with fewer measurements compared to traditional compressed sensing methods.
\\\\
A theoretical justification is given in \cite{bora2017}. If \( G \) is \( L \)-Lipschitz, a number of random Gaussian measurements of the order of \( k \log L \) is sufficient to guarantee that the reconstruction error is close to the smallest error achievable by any vector in the range of \( G \). The number of measurements required therefore scales with the latent dimension \( k \) and not with the ambient dimension \( n \), which is the reason why so few measurements are enough.
\\
In order to improve convergence we performed a thousand gradient steps for each of ten random restarts, and we kept the attempt that presented the smallest measurement error \( \| AG(z) - y \| \) at the end of the approximation. Note that the regularization term is excluded from this comparison: ranking the restarts on the full objective would reward a small \( \|z\| \) rather than a faithful reconstruction. The measurement error is also the only criterion available in practice, since it can be computed without knowing \( x^* \). These steps are not training epochs: the weights of \( G \) are frozen and the only variable that gets updated is the latent vector \( z \).

